\documentclass[conference]{IEEEtran}
\IEEEoverridecommandlockouts
\usepackage{cite}
\usepackage{amsmath,amssymb,amsfonts}
\usepackage{algorithmic}
\usepackage{graphicx}
\usepackage{textcomp}
\usepackage{xcolor}

\begin{document}

\title{Military Object Detection: A Comparative Analysis of Faster R-CNN and YOLO}

\author{\IEEEauthorblockN{1\textsuperscript{st} Group Member}
\IEEEauthorblockA{\textit{Department of Computer Science} \\
\textit{University}\\
City, Country \\
email@domain.com}
}

\maketitle

\begin{abstract}
The detection and classification of military assets such as tanks and aircraft in diverse operational environments is a critical challenge in modern surveillance. This report presents a comprehensive comparison between a two-stage detector (Faster R-CNN) and a single-stage detector (YOLO) applied to a custom military object dataset. We detail the exploratory data analysis, specifically addressing center-bias through novel spatial heatmaps, and discuss the architectural nuances, training methodologies, and deployment strategies utilizing FastAPI and Docker. Extensive evaluations demonstrate the trade-offs between precision, characterized through a background-aware confusion matrix, and real-time inference speed.
\end{abstract}

\begin{IEEEkeywords}
Object Detection, Deep Learning, Faster R-CNN, YOLO, Military Surveillance, Convolutional Neural Networks
\end{IEEEkeywords}

\section{Introduction}
Modern military operations increasingly rely on automated surveillance systems to detect and classify adversarial assets. The deployment of Convolutional Neural Networks (CNNs) has revolutionized object detection, yet applying these models to military contexts introduces unique challenges, including severe background clutter and varied object scales. In this study, we refactor an experimental object detection codebase into a production-ready pipeline, comparing the efficacy of Faster R-CNN and YOLO. 

\section{Related Work}
Object detection methodologies broadly bifurcate into two-stage approaches, pioneered by the R-CNN family, and single-stage approaches like YOLO and SSD. Two-stage models prioritize localization precision by utilizing Region Proposal Networks (RPNs) before classification. Conversely, YOLO models reframe detection as a single regression problem, optimizing for inference speed. Recent advancements have sought to harmonize these paradigms for defense applications, balancing computational efficiency with high fidelity accuracy.

\section{Methodology}
Our approach encompasses robust data preprocessing, modular architecture design, and systematic training.

\subsection{Exploratory Data Analysis}
To ensure robust model generalization, we analyzed the spatial distribution of ground-truth bounding boxes. A custom Bounding Box Spatial Heatmap was engineered to visually prove the absence of center-bias within the dataset, ensuring the models learn intrinsic object features rather than relying on location priors.

\subsection{Faster R-CNN Implementation}
The Faster R-CNN model was instantiated with a ResNet-50 backbone coupled with a Feature Pyramid Network (FPN). The training loop was isolated utilizing modular PyTorch data loaders. A critical enhancement included the implementation of a custom confusion matrix that explicitly models the 'Background' class, allowing for rigorous analysis of false positives generated by environmental clutter.

\subsection{YOLO Implementation}
The YOLO architecture was configured to maximize throughput for potential edge deployment. We instituted a Qualitative Failure Analysis Grid, which systematically sorts validation predictions by Confidence and Intersection over Union (IoU) to visualize the top False Positives and False Negatives, thereby identifying systemic model weaknesses.

\section{Experiments and Results}
The models were evaluated on a held-out test set utilizing standard COCO metrics (mAP@0.5:0.95). 

\subsection{Performance Metrics}
Faster R-CNN achieved superior average precision, particularly for small, occluded objects, corroborating its robust region proposal mechanism. Conversely, YOLO exhibited unparalleled inference speeds, making it viable for real-time video analytics.

\subsection{Deployment}
A robust deployment pipeline was constructed using FastAPI to serve the optimal model via a RESTful endpoint. The application is containerized within an NVIDIA GPU-optimized Docker environment, ensuring scalable and reproducible inference capabilities.

\section{Conclusion}
This project successfully transformed an experimental military object detection framework into a professional, scalable repository. We empirically demonstrated the trade-offs between the precision of Faster R-CNN and the speed of YOLO. Future work will explore integrating Transformer-based architectures and employing synthetic data generation to further augment the training corpus.

\end{document}
